\cotefile{34}
\scanpage{34}
\mspage{32}
% archive sheet number: 33
% image line 1
\deleted{Ce \unsure{cas|cor.}, \unsure{est évident} grâce à \unsure{8.6|6.6}.}

7. \uline{Le théorème de monodromie :} \deleted{\unsure{démonstration arithmétique}} \inserted{\uline{Énoncé conjectural}.}

\deleted{L'énoncé est \unsure{le suivant}.}

\uline{Conjecture} 7.1. \inserted{(Serre-Tate)} \unsure{Sous|Soit|Soient}
\deleted{$X$ un schéma de type fini sur $K$, $i$ un entier, $\ell$ un nb premier $\ne p$, \unsure{et $i$ un ent.} \unsure{On suppose} que $H^i(X_{\overline K},\unsure{\mathbb Z/\ell\mathbb Z|Z/\ell Z})$ est fini. Cf. 0.3)}
les conditions de l'énoncé 0.3, si on suppose $\ell\ne p$, la représentation induite sur le groupe d'inertie
\[
 I\longrightarrow \operatorname{Aut}_{\mathbb Z_\ell}\bigl(H^i(X_{\overline K},\mathbb Z_\ell)\bigr)
\]
est quasi-unipotente, i.e. $\exists$ s-groupe d'indice fini $U$ dans $I$ tel que l'image de $U$ est formée d'\unsure{opérations|opérateurs} unipotentes dans $H^i(X_{\overline K}\,\mathbb Z_\ell)$.

7.2. Nous allons voir par la suite deux démonstrations ``relatives'' de cette conjecture, utilisant l'une et l'autre la résolution des singularités, \deleted{\unsure{et les th de pureté}} de façon d'ailleurs assez différente. De ces démonstrations, il résultera que la conjecture est vraie dans chacun des cas suivants :

a) $p=1$, i.e. $k$ (ou encore $S$) de car. nulle.

\margin{\deleted{\unsure{$X$ est propre}}}
b) $i=1$,

c) $S$ est le hensélisé \inserted{(ou le hensélisé strict)} d'un anneau local d'un point d'un schéma de type fini sur un corps, ou sur $\operatorname{Spec}.\mathbb Z$, et $X_K$ est propre sur $K$ \uline{ou} $K$ de car. nulle.

La première démonstration, qui est de nature arithmétique, sera donnée dans le
\scanpage{35}
\mspage{33}
% archive sheet number: 34
présent exposé en se réduisant cf. \unsure{6.6|8.6}. à \deleted{\unsure{un cas}} une situation avec un corps résiduel fini. \inserted{Cette démonstration \deleted{\unsure{s'applique pour}} donne la conj. dans les cas a) et c).}
\deleted{Dans} La démonstration géométrique fera l'objet de l'exposé suivant; elle permettra de préciser assez considérablement \deleted{l'énoncé et} la conclusion de unipotence, et \deleted{\unsure{permet de}} s'applique à l'heure actuelle dans les cas a) et b).

7.3. On peut donner une forme précisée de 7.1, en y faisant varier $\ell$ parmi tous les nbs premiers distincts de $p$. Serre et Tate \deleted{se bornent à} demandent que les \deleted{représentations obtenues} groupes $\Gamma$ \deleted{\unsure{opérant dans}} (quotients de $I$) associés à \deleted{\unsure{deux repr. distinctes}} \inserted{$H^i$ pour deux $\ell$ distincts} soient les mêmes, et que les représentations obtenues de $\Gamma$ soient ``isomorphes'' (sur une \unsure{sur-extension} commune $\unsure{\Omega}$ de $\unsure{\mathbb Q_\ell|Q_\ell},\mathbb Q_{\ell'}$); \deleted{\unsure{et}} \inserted{Plus précisément, on peut demander} \deleted{\unsure{la relation}} \inserted{$\operatorname{Tr}g_{H^i(\overline X,\mathbb Q_\ell)}$ soit un entier naturel pour $\forall\,g\in I$, et qu'il soit indépendant de $\ell\ne p$:}
\deleted{\unsure{s'exprime}} \unsure{plus simplement}
\[
 \operatorname{Tr}g_{H^i(\overline X,\mathbb Q_\ell)}
 =\operatorname{Tr}g_{H^i(\overline X,\mathbb Q_{\ell'})}
 \tag*{(7.3.1)}
\]
pour tout $g\in\pi$. [Ceci sera prouvé pour $i=1$ par la suite, en utilisant les modèles de Néron]. \deleted{\unsure{De plus, on est conduit à une question plus considérable}.}

7.4. Lorsque $k$ est fini, \deleted{Serre et Tate} il y a lieu (avec Serre et Tate) de conjecturer que cette \deleted{égalité} \inserted{relation (7.3.1)} est vraie \inserted{(ds $\mathbb Q$)} pour $\forall\,g\in\pi$ dont l'image dans $\operatorname{Gal}(\overline k/k)\simeq\widehat{\mathbb Z}$ tombe dans $\mathbb Z$. [Cela même \deleted{\unsure{pour}} \inserted{sera prouvé} pour $i=1$ dans notre séminaire.]

On veut de plus que si $g\in\pi$ se \unsure{projette} sur \deleted{\unsure{f}} \inserted{$f^{-1}$} (f le générateur
% margin note written slanted
\margin{ou frobenius) les \uuline{valeurs} propres de $g$ opérant sur $H^i(\overline X,\mathbb Q_\ell)$ soient de \unsure{valeur|valeurs} \unsure{absolue complexe|absolues complexes} de la forme $q^{\unsure{w|u}/2}$, où $q=\operatorname{card}k$, $0\leq w\leq2i$.}

7.4. On peut préciser partiellement ces conjectures, en \unsure{tenant|donnant} compte des représentations de tout le groupe algébrique $\mathbb G_{a\,\mathbb Q_\ell}\times\Gamma_{\mathbb Q_\ell}$, pas seulement de $\Gamma_{\mathbb Q_\ell}$,
\scanpage{36}
\mspage{34}
% archive sheet number: 35
en remarquant que ce groupe algébrique sur $\unsure{\mathbb Q_\ell|Q_\ell}$ provient d'un groupe algébrique \deleted{sur} $\mathbb G_{a\,\mathbb Q}\times\Gamma_{\mathbb Q}$ sur $\mathbb Q$. On peut demander \inserted{d'abord} \deleted{que} si les représentations algébriques sur les différents $\mathbb Q_\ell$ définies par les $H^i(\overline X,\mathbb Q_\ell)$ proviennent d'une même représentation algébrique \uline{sur} $\mathbb Q$.
\deleted{(Ça devrait \unsure{sortir|rester} \unsure{un exercice|en exercice} \unsure{pour $i=1$}.)}
\inserted{Ce n'est raisonnable \deleted{que} cependant qu'en car. nulle, où une démonstration par voie transcendante (utilisant la cohomologie entière transcendante) devrait être possible. En car. résiduelle $>0$, on peut \unsure{vouloir|voir} que ces représentations s'écrivent (de façon unique à isom près) sur $\mathbb Q$; on peut demander qu'elles y soient indépendantes de $\ell$.}
% slanted
\margin{[N.B. \unsure{Même|M\textsuperscript{me}} la représentation induite de $\Gamma$ ne s'écrit pas en général sur $\mathbb Q$, si $i=1$ et si $X_K$ est une courbe elliptique sur $K$ ayant ``bonne réduction potentielle''.]}

Dans le cas d'un corps résiduel fini, on trouve un groupe ``loc. algébrique''
\[
 \mathbb Z_{\mathbb Q}\cdot(\mathbb G_{a\,\mathbb Q}\times\Gamma_{\mathbb Q})=G
\]
sur $\mathbb Q$, produit semi-direct, où \inserted{le générateur} $f_q\in\mathbb Z_{\mathbb Q}(\mathbb Q)$ opère sur $\mathbb G_{a\,\mathbb Q}$ \deleted{\unsure{par une}} par multiplication par $q$, et sur $\Gamma$ comme il faut [sur $\unsure{(\Gamma/[\Gamma,\Gamma])(\ell)|(\Gamma/[\Gamma,\Gamma])^{(\ell)}|(\Gamma/[\Gamma,\Gamma])_\ell}$ il opère aussi par multiplication par $q$]. On peut encore demander que les représentations sur les différents $\mathbb Q_\ell$ ($\ell\ne p$) définies par les $H^i(\overline X,\mathbb Q_\ell)$ \deleted{soient définies sur $\mathbb Q$, et déduites} \inserted{isomorphes à celles déduites} d'une même représentation de $G$, sur $\mathbb Q$.
\scanpage{37}
\mspage{35}
% archive sheet number: 36
8. \uline{Théorème de monodromie : démonstration arithmétique}

\uline{Th.} 8.1. Supposons que $V$ soit limite inductive \inserted{filtrante} de sous-anneaux locaux $V_i$, \deleted{l'idéal} \inserted{d'idéal maximal $\unsure{\mathfrak m_i|m_i}=\mathfrak m\cap V_i$} contenant l'uniformisante $t$, tels que les deux conditions suivantes soient satisfaites :

a) \deleted{\unsure{le div}} $V_i$ \unsure{est} \inserted{et $V_i/tV_i$ sont} réguliers, \deleted{et le diviseur $\operatorname{div}(t)$ \unsure{est} régulier, i.e. $t$ fait partie d'un système de param. de $V_i$.} \deleted{dans $\operatorname{Spec}V_i$, est à croisements normaux.}

b) Le corps \inserted{résiduel} $k_i=k(V_i)$ est tel que aucune extension finie de $k_i$ ne contienne toutes les racines $\ell^{\unsure{\nu|n}}$-ièmes de l'unité, pour $\forall$ \unsure{$\ell$ entier $>1$}.

Alors le th. de monodromie \unsure{sous} la forme 7.1. est vrai pour des schémas \deleted{de t.f.} \inserted{\uline{\unsure{propre|propres}}} $X$ sur $K$, \deleted{tel que $\unsure{H^i|H^0}(X_{\overline K},\unsure{\mathbb Z/\ell\mathbb Z|Z/\ell Z})$ est fini} \inserted{lorsque l'opération de $I$ sur $H^i(X_{\overline K},\mathbb Z_\ell)\otimes_{\mathbb Z_\ell}\unsure{\mathbb Q_\ell|Q_\ell}$ se fait à travers $I_t$,} \deleted{dans} \uline{ou} que les $V_\alpha$ sont des anneaux de val. discrète.

\uline{Corollaire} 8.2. Le théorème de monodromie \inserted{sous la forme 7.1.} est vrai \deleted{lorsque $H^i(X_{\overline K},\mathbb Z/\ell\mathbb Z)$ est fini,} \inserted{pour $X_K$ propre sur $K$} dans chacun des cas suivants :

a) $V$ est \inserted{le hensélisé ou le hensélisé strict de} l'anneau local d'un point d'un schéma $S$ de type fini sur un corps $k$, ou sur \deleted{\unsure{Spec}} \inserted{un localisé de} $\mathbb Z$.

b) $V$ est de car. nulle, (i.e. \inserted{$k$ est} de caractéristique \deleted{\unsure{résiduelle}} nulle).

\deleted{Cas a)}
\unsure{Démonstration}. 8.2.

\uline{Cas} a). \deleted{Supposons $V$ le hensélisé de l'anneau local d'un \unsure{point}. \unsure{Si} $S$ est de type fini sur $\operatorname{Spec}\mathbb Z$,}
\unsure{Supposons} d'abord que l'anneau de val. \deleted{\unsure{discr.}} $V$ soit l'anneau local d'un schéma $S$ de type fini sur \deleted{\unsure{Spec}} \inserted{un localisé} $\mathbb Z$, ou sur un corps $k_0$. Dans le premier cas, \inserted{ou si $k_0$ est de type fini sur le corps premier,} le corps résiduel $k(V)$ est de type fini sur le corps premier, donc $V$ satisfait la condition
\scanpage{38}
\mspage{36}
% archive sheet number: 37
Dans le cas d'un corps de base $k_0$ quelconque, $k$ \deleted{\unsure{$S$ est li}} est limite inductive de ses sous-extensions $k_\alpha$ de t.f. sur le corps premier, et pour $\alpha$ assez grand $\unsure{S_0|S}$ provient d'un $S_\alpha$ sur $k_\alpha$. Soit $x_\alpha$ l'image de $x\in S$ dans $S_\alpha$. Comme $\unsure{\mathscr O_{S,x}|\mathcal O_{S,x}}$ est \uline{plat} sur $\mathscr O_{S_\alpha,x_\alpha}$, on a $\dim\mathscr O_{S_\alpha,x_\alpha}\leq\dim\mathscr O_{S,x}=1$, \uline{et} $\mathscr O_{S_\alpha,x_\alpha}$ est régulier, donc $\mathscr O_{S_\alpha,x_\alpha}$ est un corps ou \deleted{\unsure{un anneau de}} un anneau de val. discrète.
D'ailleurs pour $\alpha$ assez grand, c'est un anneau de valuation discrète, \deleted{et si $\unsure{\mathfrak m_\alpha|m_\alpha}$ est} (sinon $\mathscr O_{S,s}$ serait un corps!), \deleted{de plus $\mathfrak m_\alpha$ est engendré par $t$} et $t\in\mathfrak m_\alpha$. \deleted{Comme \unsure{$\mathfrak m/\mathfrak m^2=\lim\mathfrak m_\alpha/\mathfrak m_\alpha^2$}, on voit que} $t\notin\mathfrak m_\alpha^2$ (sinon on aurait $t\in\mathfrak m^2$), donc on est sous les conditions de a), b) de 8.1.
\scanpage{39}
\mspage{37}
% archive sheet number: 38
\deleted{de 8.1. Dans le \unsure{deuxième cas}, \inserted{cas d'un corps quelconque} \unsure{pour des généralités} \unsure{dans une} $V\otimes_{k_0}k'_0$, où $k'_0$ est une extension \unsure{radicielle|radicale} convenable de $k_0$, $k_0=\underset{\longrightarrow}{\lim}\,k_\alpha$, où $k_\alpha$ parcourt les sous-corps de $k_0$ qui sont des extensions de type fini \unsure{dans le} corps premier. Pour $\alpha$ assez grand, $S$ provient de $S_\alpha$ \inserted{lisse} sur $k_\alpha$, \inserted{et le diviseur $(t)$ de $S$ provient d'un diviseur \unsure{réduit} $D_\alpha$ sur $X_\alpha$.} Alors $V$ est limite des $\unsure{\mathscr O_{S_\alpha,x_\alpha}|\mathcal O_{S_\alpha,x_\alpha}}$, $x_\alpha$ étant \unsure{le point considéré dans $S_\alpha$} \inserted{le point générique de $D_\alpha$}. Or le \unsure{corps} $\mathscr O_{\unsure{S_\alpha|S_0},x_\alpha}$ est un anneau de valuation discrète à corps résiduel de type fini (au sens absolu).}
% slanted
\margin{\deleted{(ce qui ne change pas les groupes d'inertie), on est ramené au cas où $k(V)$ sur $k$ est \uline{séparable}. On peut donc supposer $S$ \uline{lisse} sur $k_0$, et \unsure{Écrivons}.}}

Dans \unsure{donc} les deux cas, l'anneau \deleted{\unsure{ainsi}} \inserted{(resp. $V\otimes_{k_0}k'_0$)} est \uline{lim} filtrante \inserted{(par morphismes d'inclusions locaux)} de \deleted{ses} sous-anneaux de valuation discrets, \inserted{$V_\alpha$ à corps résiduel de type fini. Donc,} \deleted{\unsure{avec corps rés. de type fini de $k$ considéré}.} le hensélisé $\unsure{V^h|V^{sh}}$ (resp. $(V\otimes_{k_0}k'_0)^h=V^h\otimes_{k_0}k'_0$) est limite inductive des $V_\alpha^h$, et on peut lui appliquer 8.1. De même pour le hensélisé strict. \deleted{\unsure{8}}

b) Cas $V$ de car. nulle, \inserted{i.e. algèbre sur $\unsure{\mathbb Q|Q}$.} $V$ est limite inductive de \inserted{ses} sous-$\mathbb Q$-algèbres locales $V_\alpha$, ess. de type fini sur $\mathbb Q$, \deleted{avec $V_\alpha$} d'idéaux maximaux $\unsure{\mathfrak m|m}\cap V_\alpha=\mathfrak m_\alpha$. On peut supposer $t\in V_\alpha$. Par Hironaka, on peut trouver un \deleted{\unsure{modèle normal}} schéma propre, \deleted{sur} intègre, \uuline{régulier} $S'_\alpha$ sur $S_\alpha=\operatorname{Spec}V_\alpha$, avec $S'_\alpha\to S_\alpha$ birationnel, et l'image inverse \inserted{dans $S'_\alpha$} \deleted{est} de $\operatorname{div}(t_{S_\alpha})$ un diviseur à croisements normaux. \deleted{\unsure{Prenant}} Comme $S'_\alpha-S_\alpha$ est propre, il existe un point $x'_\alpha$ de $S'_\alpha$ tel que l'anneau local $V'_\alpha\unsure{:|=}\mathscr O_{\unsure{S'_\alpha|\alpha'},x'_\alpha}$ soit dominé par $V$. \deleted{Donc} Or cet anneau local est régulier, et le diviseur de $t$ y est à croisements normaux.
% slanted
\margin{\deleted{\unsure{L'hypothèse}} donc $t=\prod x_i^{n_i}$, \unsure{les $x_i\in\mathfrak m'_\alpha$|les $x_i\in\mathfrak m_\alpha$}, $n_i\geq1$. D'où le fait qu'il n'y a qu'un seul $x_i$, et que \deleted{\unsure{de plus on a}} $\unsure{n_1|n_i}=1$.}
\deleted{Donc} les $V_\alpha$ ayant cette propriété supplémentaire
\scanpage{40}
\mspage{38}
% archive sheet number: 39
($V_\alpha$ régulier, \deleted{$\operatorname{div}(t_{V_\alpha})$ rég.} $V_\alpha/tV_\alpha$ régulier) sont \uline{cofinaux} parmi les $V_\alpha$. De plus, tout $V_\alpha$ a un corps résiduel de type fini. On est donc \unsure{sous|dans} les conditions de 8.1.

\uline{Remarque} 8.3. \deleted{\unsure{Cette dé}} La démonstration du cas b) \deleted{serait} serait valable sans restriction sur $V$, une fois qu'on disposerait de la résolution des singularités \inserted{sous la forme} de Hironaka \deleted{\unsure{sur le}} (avec diviseurs à croisements normaux $\ldots$) pour les schémas de \uline{type fini sur $\unsure{\mathbb Z|Z}$}; \inserted{cela impliquerait que l'hyp. de 8.1. sur $V$ est \uline{\unsure{tjs|tjrs}} satisfaite.} Elle marcherait de même si $V$ est de car. $p>0$, si on dispose de la résolution des singularités à la Hironaka pour les schémas de t.f. sur $\mathbb F_p$. Notons qu'on n'a besoin ici \deleted{\unsure{de la de Hi}} que de ``l'\uline{uniformisation locale}''.

8.4. \uline{Démonstration de} 8.1. On \deleted{\unsure{aura}} \inserted{peut supposer les $V_\alpha$ henséliens,} quitte à les remplacer par leurs hensélisés. On aura
\[
 K=\underset{\longrightarrow}{\lim}\,\underbrace{(V_\alpha)_t}_{K_\alpha},
\]
donc pour $\alpha$ assez grand, $X_K$ provient d'un schéma \deleted{de type f.} $X_{\deleted{\unsure{K}}\alpha}$ sur \deleted{\unsure{Spec}} $(V_\alpha)_t$, i.e. sur l'ouvert $U_\alpha=S_{\alpha t}$ de $S_\alpha=\operatorname{Spec}V_\alpha$; \deleted{\unsure{prenons}} soit $\eta_\alpha$ le pt générique de $S_\alpha$ i.e. de $S_{\alpha t}$, c'est l'image de $\eta$.
\deleted{Comme $f_\alpha:X_{K_\alpha}\to U_\alpha$ est \unsure{Choisissons un ouvert},}
Comme $f_\alpha$ est propre, on sait (SGA 5 \uline{VI}) que $\unsure{E_{U_\alpha}|\mathscr E_{U_\alpha}}=R^if_{\alpha *}(\mathbb Z_\ell)$ sur $U_\alpha$ est constructible, donc \deleted{$\exists$ \unsure{un}} voisinage ouvert $U'_\alpha$ de $\eta_\alpha$ sur lequel $E_\alpha$ est \unsure{constant tordu}. Augmentant au besoin $\alpha$, \inserted{comme $g^{-1}(U_\alpha)=\unsure{g|j}^{-1}(\unsure{V_\alpha|U'_\alpha})=\{\eta\}$,} on
\diagram{\begin{tikzcd}[ampersand replacement=\&]
S \arrow[d,"g"'] \& \eta_\alpha \\
S_\alpha\supset U_\alpha \&
\end{tikzcd}}
\scanpage{41}
\mspage{39}
% archive sheet number: 40
peut supposer $\unsure{U'_\alpha|V_\alpha}=U_\alpha$ (EGA \uline{IV} 8), donc $\unsure{E_\alpha|\mathscr E_\alpha}$ est un faisceau constant tordu sur $U_\alpha$, \inserted{donc provient d'une représentation $\ell$-adique de $\pi_\alpha=\pi_1(U_\alpha,\overline\eta)$ dans $\unsure{\operatorname{Aut}E}$.} En vertu du th. de changement de base propre (SGA 4 \uline{XII}), $Rf_{K*}((\unsure{\mathbb Z_\ell|Z_\ell})_{X_K})=E_K$ est \deleted{\unsure{lui}} isomorphe à l'image inverse de $E_{U_\alpha}$ par \deleted{\unsure{de $g$}} $\eta\longrightarrow U_\alpha$, donc la représentation de $\pi$ \deleted{\unsure{provenant}} définie par $H^i(X_{\overline K},\mathbb Z_\ell)$ est celle déduite de $\pi_\alpha$ par l'hom canonique $\pi-\pi_\alpha$.
\margin{\unsure{donc est} $E$ de t.f. sur $\mathbb Z_\ell$}
\[
\deleted{\pi\longrightarrow\pi_\alpha}.
\]
\deleted{Notons que $\pi_\alpha$ est une extension de $\pi_{\alpha,0}=\operatorname{Gal}(\overline k_\alpha/k_\alpha)$}

Pour la suite, mettons l'exposant $'$ au lieu de l'indice $\alpha$, et écrivons
\[
\unsure{\varphi|\alpha}\qquad \pi-\pi'.
\]
Notons que $\pi'$ est une extension de
\[
\pi'_0=\operatorname{Gal}(\overline{k'}/k')\simeq\pi_1(S'_{\deleted{\unsure{\alpha}}})
\]
par un ss-groupe \deleted{\unsure{de}} $I'$, et que l'hom $\varphi$ induit un hom :
\[
I\longrightarrow I'\subset\pi'.
\]
Dans le cas où $V_\alpha$ est de val. discrète, \deleted{\unsure{par}} la conclusion résulte de n° 6.
% slanted
\margin{Ceci \unsure{suffit|suffirait} déjà à entraîner le cas 8.2 a), qui ne sera pas \unsure{démontré} pour l'instant par la dém. ``géom.'' de l'exposé suivant (\unsure{ce nom est|ceci étant} \unsure{la seule application|la seule appellation indiquée} pour le moment de la méthode ``arithmétique'', \unsure{le lecteur pourra|la lecture permettant} \unsure{se dispenser de lire|de dispenser le lecteur, si} \unsure{le cas} \uuline{général} \unsure{le gêne} !)}

\deleted{Nous allons utiliser donc le théorème de structure suivant, qui est une \inserted{conséquence} \unsure{variante} du ``lemme de Abhyankar'' :}

\deleted{\uline{Lemme} 8.5. Soit $A$ un anneau local noethérien \uline{hensélien}, $k$ le corps résiduel \unsure{$k$}, $p$ l'exp. caractéristique de $k$, \unsure{Spec} $S=\operatorname{Spec}A$, $D=\sum_{i\in I}D_i$ un diviseur à croisements normaux, $(x_i)_{i\in I}$ une partie d'un syst. de paramètres, $a=\prod x_i$, $U=S_a$, $\unsure{\xi|\zeta}$ un pt géom. de $U$. Alors}
\deleted{Alors $\xi$ définit une \deleted{extension séparable} \inserted{clôture séparable} $\overline k$ de $k$, et on}
\scanpage{42}
\mspage{40}
% archive sheet number: 41
\deleted{a un isom. canonique}
\[
\deleted{\pi_1(S,\unsure{\xi|\zeta})\simeq\pi_1(\operatorname{Spec}k,\overline k)=\operatorname{Gal}(\overline k/k)}.
\]
\deleted{De plus, l'hom. canonique (cf.)}
\[
\deleted{\pi_1(U,\xi)\longrightarrow\pi_1(S,\xi)=\pi_0}
\]
\deleted{est surjectif, \deleted{et définit une extension dont le noyau} $I$ a la structure suivante :}

\deleted{Soit} $\unsure{E|\mathscr E}$

Quitte à diviser $E$ par son s-groupe de torsion, \inserted{(ce qui ne change pas la conclusion à prouver, ni l'hyp.)} on est ramené à examiner le cas d'un $E$ libre sur $\unsure{\mathbb Z_\ell|Z_\ell}$, l'hyp. d'action modérée s'exprimant maintenant par le fait que l'image de $I$ dans $\operatorname{Aut}(E_0)$ ($E_0=E\otimes_{\mathbb Z_\ell}\mathbb Z/\ell\mathbb Z$) est d'ordre premier à $p$. Mais on voit de suite que $I$ est la limite projective des groupes $I_\alpha$, donc l'image de $I$ dans $\operatorname{Aut}(E_0)$ est \deleted{\unsure{lim}} l'intersection de celles des $I_\alpha$, donc est égale à celle de $I_\alpha$ pour $\alpha$ grand. Donc on peut supposer \deleted{que} \inserted{$\alpha$ tel que} cette image est d'ordre premier à $p$. Donc $I'=I_\alpha$ opère \inserted{sur $E$} par l'intermédiaire d'un groupe profini \unsure{pro-fini|pro-d'ordre} de l'\deleted{entier} \inserted{ordre} transfini $\unsure{q}$ premier à $p$.

\unsure{b} On utilise maintenant le th. de structure suivant, qui est une des variantes du lemme de Abhyankar :
% slanted
\margin{devrait faire l'objet d'un n° spécial avant l'actuel n° 2}

\uline{Lemme} 8.5. Soit $A$ anneau local noeth. régulier \uline{hensélien}, $k$ son corps résiduel, $p$ l'exp. car. de $k$, \deleted{$\unsure{(x_i)_{i\in I}}$} $(x_i)_{1\leq i\leq n}$ une partie d'un syst \uline{régulier}
\scanpage{43}
\mspage{41}
% archive sheet number: 42
de paramètres de $A$, $x=\prod x_i$, \deleted{$\unsure{U=}$} $S=\operatorname{Spec}A$, $U=\operatorname{Spec}A_x=S_x$, $\unsure{\xi|\zeta}$ un pt géom. de $U$.

a) $\xi$ définit une clôture séparable $\overline k$ de $k$ et on a un isom canonique
\[
\pi_1(S,\xi)\simeq\pi_1(\operatorname{Spec}k,\overline k)=\operatorname{Gal}(\overline k/k)
\tag*{(8.5.1)}
\]

b) L'hom canonique
\[
\pi_1(U,\xi)\longrightarrow\pi_1(S,\xi)\simeq\operatorname{Gal}(\overline k/k)
\tag*{(8.5.2)}
\]
est surjectif, et définit donc $\pi_1(U,\xi)$ comme extension de $\operatorname{Gal}(\overline k/k)$ par un groupe \deleted{$\unsure{I}$} profini $I$; \deleted{\unsure{qui est} \unsure{canoniquement} \unsure{isomorphe} à $\pi_1(\widetilde S,\widetilde\xi)$, où $\widetilde S$ est le hensélisé strict}
\[
1-I-\pi_1(U,\xi)-\operatorname{Gal}(\overline k/k)-1,
\tag*{(8.5.3)}
\]
et on a un isom. can.
\[
I\simeq\pi_1(\widetilde U,\widetilde\xi),
\tag*{(8.5.4)}
\]
où $\widetilde S$ désigne le hensélisé strict de $S$ relativement à $\overline k/k$, \deleted{\unsure{et où}} $\widetilde U$ l'image inverse de $U$ dans $\widetilde S$, et $\widetilde\xi$ un relèvement canonique de $\xi$ à $\widetilde S$.

c) Soit $I_t$ le plus grand \unsure{pro-quotient|groupe quotient} de $I$ qui est de cardinal transfini premier à $p$. Alors on a un isom canonique
\[
I_t\simeq\deleted{[\widehat{\unsure{\mathbb Z|Z}}'(1)]^n}\left(\prod_{\ell\ne p}\mathbb Z_\ell(1)\right)^n;
\tag*{(8.5.5)}
\]
\[
\mathbb Z_\ell(1)\simeq T_\ell(k(\xi)^*)\simeq T_\ell(\overline k^*).
\tag*{\uline{(8.5.6)}}
\]
\deleted{\unsure{voir|où l'on a}}
\[
\deleted{\widehat{\mathbb Z}'(1)=\prod_{\ell\ne p}\mathbb Z_\ell(1)}
\]
% slanted
\margin{ne sert pas pour la dém. de 8.1.}
(Pour un revêtement principal donné $S'$ de $S$, de groupe de Galois $\Gamma$, la représentation correspondante $\pi\longrightarrow\Gamma$ induit un hom $I\longrightarrow\deleted{\unsure{I}}\Gamma$; que ce dernier se factorise par $I_t$ \uline{si et seulement si} \deleted{soit}
\scanpage{44}
\mspage{42}
% archive sheet number: 43
pour $\forall\,1\leq i\leq n$, $S'$ est modérément ramifié en les pts \deleted{\unsure{résiduels du diviseur}} \inserted{génériques de $\operatorname{div}(x_i)$}.

d) L'action naturelle de $\pi_1(U,\unsure{\xi|\zeta})$ sur $\inserted{I_t=}\prod_{\ell\ne p}I/[I,I]^{(\ell)}$, qui se factorise \deleted{\unsure{par}} évidemment par une action de $\operatorname{Gal}(\overline k/k)$, s'identifie via l'isomorphisme (8.5.5) à l'action \deleted{de $\operatorname{Gal}(\overline k/k)$} déduite de l'action naturelle de $\operatorname{Gal}(\overline k/k)$ sur $T_\ell(\overline k^{\unsure{\times|*}})$, \deleted{\unsure{déduite}} via les opérations de $\operatorname{Gal}(\overline k/k)$ sur les ${}_{\ell^\nu}(\overline k^*)\subset\overline k^*$.

Démonstration laissée au lecteur \unsure{---}.

\uline{Corollaire} 8.6. Sous les conditions de 8.5, supposons qu'aucune extension finie de $k$ ne contienne toutes les racines $\ell^\nu$-ièmes de $1$, pour $\forall$ entier $\nu$. Alors, pour $\forall$ représentation $\ell$-adique de $\pi_1(U,\xi)$ \inserted{``modérée'', i.e. telle que} l'action induite de $I$ \deleted{est quasi-unipotente} \inserted{se factorise par une action de $I_t$, cette dernière est} \uline{quasi-unipotente}.

\uline{Remarque 8.7} \unsure{a)} Dans le cas d'application de 8.6, on aura
\deleted{$t=u\prod x_i^{\alpha_i}$, avec des $\alpha_i\geq1$, et $x_i\in\unsure{\mathfrak m_i|m_i}\subset\mathfrak m=tV$, \inserted{$u_\alpha\in A^*\subset V^*$,} ce qui prouve que $n=1$ (un seul $x_i$ \inserted{soit $x_1$}) et $\alpha_1=1$. Donc \unsure{on a}}
\deleted{Alors la connaissance de la structure du groupe}
$n=1$, et l'hom $I\longrightarrow I'$ induit un \uline{isom} $I_t\xrightarrow{\sim}I'_t$.

\uline{Remarque} \deleted{\unsure{b}} 8.8. La démonstration de 8.1. marche sans hyp. de propreté sur $X_K$, si on sait que l'on peut trouver un voisinage ouvert $U_\alpha$ du pt générique $\eta_\alpha$ de $S_\alpha$
\scanpage{45}
\mspage{43}
% archive sheet number: 44
au dessus duquel $\inserted{\unsure{E_\alpha|\mathscr E_\alpha}=}R^if_{\alpha *}(\unsure{\mathbb Z_{\ell\,X_\alpha}|Z_{\ell\,X_\alpha}})$ est constant tordu. En effet, le fait que \deleted{$\unsure{E}$} $R^if_{K*}(\mathbb Z_{\ell\,X_K})$ s'en déduise par le chgt de base $\eta\longrightarrow\eta_\alpha\longrightarrow S_\alpha$ signifie que la formation de $R^if_{\alpha *}(\mathbb Z_{\ell\,X_\alpha})$ commute avec les chgts de base $\eta_\alpha-S_\alpha$ et $\eta-\eta_\alpha$. Pour le premier c'est trivial, pour le deuxième c'est une conséquence du th. de chgt de base \deleted{lisse} par morphisme lisse (SGA 4 \uline{XVI} 1.6). \deleted{D'autre part, l'exis}
D'autre part, l'existence d'un $\unsure{V_\alpha|U_\alpha}$ comme dessus est \unsure{encore prouvée|encore garantie}, \unsure{disposant de|si on dispose de} la résolution des singularités \inserted{(sous la forme de Hironaka)} pour les schémas de type fini et de dimension $\leq\dim X_{K\deleted{\unsure{\alpha}}}$ sur $K_\alpha$ (SGA 5 \uline{II} 5.1). \deleted{\unsure{On le voit ainsi}} \inserted{Cette hypothèse est vérifiée} en particulier, par Hironaka, si $K$ est de car. nulle, et \deleted{\unsure{justifie le th. de monodromie lorsque les termes dans}} \inserted{par Abhyankar si $\dim X_K\leq2$.}

\deleted{\uline{Corollaire} 8.9. Soit $V$ comme dans 8.1, \inserted{$X$ de type fini sur $X_K$.} \deleted{\unsure{sur $K=k(V)$}}, \deleted{\unsure{de caractéristique}}, \inserted{\unsure{hens.} i.e. $\mathbb Z$,} \uline{caractéristique} \inserted{\uline{nulle}, ou $\dim X_K\leq2$}, \deleted{soit $X_K$ de t.f. sur $K$}, \deleted{\unsure{Supp.}} \deleted{\unsure{pour}} et \unsure{supposons} les $V_\alpha$ de valuation discrète (cas de 8.2. a) p.ex) ou l'action de \deleted{\unsure{resp.}} $I$ dans $H^i(X_{\overline K},\mathbb Z_\ell)$ modérée (cas de 8.2 b) par exemple). Alors l'\deleted{\unsure{représentation}} \inserted{action} de $I$ sur $H^i(X_{\overline K},\mathbb Z_\ell)$ est quasi-unipotente.}
\margin{\deleted{\uline{Supposons $K$ de}}}

\uline{Corollaire} 8.9. Les conclusions de 8.2 restent vraies en supposant seulement $X$ de t.f. sur $K$ (pas néc. propre), pourvu que dans le cas a), on suppose la caractéristique de $K$ égale à $0$, \uline{ou} $\dim X_K\leq2$.
\scanpage{46}
\mspage{44}
% archive sheet number: 45
9. \uline{Réductions au cas projectif et lisse.}

9.0. Nous nous posons la question de savoir, pour un \inserted{$V$ et} un entier $\unsure{n|\ell}$ donné, si le th. de monodromie 7.1. est vrai pour tout $X_K$ de t.f. sur $K$, \deleted{de dim} \inserted{tel que}
\[
\dim X_K\leq n.
\]
Nous allons d'ailleurs envisager la conjecture correspondante également pour les $\pi$-modules $H^i_!(X_{\overline K},\deleted{\unsure{\mathbb Z|Z}}\unsure{\mathbb Q_\ell|Q_\ell})$ (pour $X_K$ \uline{séparé}).
\inserted{Nous supposons que la conjecture est prouvée déjà pour les dim. $\leq n-1$.}

9.1. \uline{Cas $X$ lisse.} \inserted{\uline{Relation entre $H^i$ et $H^i_!$}} Dans ce cas, \deleted{\unsure{et}} \inserted{et si $X$ est séparé,} on a un isom canonique
\[
H^i(X_{\overline K},\deleted{\unsure{\mathbb Z}}\mathbb Q_\ell)\simeq H^{\deleted{\unsure{i}}\inserted{2n-i}}_!(X_{\overline K},\mathbb Q_\ell)'\deleted{\unsure{\otimes_{\mathbb Q_\ell}\mathbb Q_\ell}}(-n),
\]
où \deleted{$\unsure{\otimes}$} $(-n)$ désigne le twist de Tate $\otimes_{\mathbb Q_\ell}\mathbb Q_\ell(-n)$, \deleted{\unsure{et}} et $'$ désigne le dual. Cet isom. est compatible avec l'action de $\pi$ et à fortiori de $I$. Comme $I$ opère trivialement sur $\mathbb Q_\ell(-n)$, il s'ensuit que l'action de $I$ sur $H^i(X_{\overline K};\mathbb Q_\ell)$ est quasi unip. (resp. unipot.) ssi son action sur $H^i_{\deleted{\unsure{*}}!}(X_{\overline K},\mathbb Q_\ell)$ l'est.

9.2. \deleted{\unsure{Le cas $X_K$ séparé}} \inserted{\uline{Cas} $H^i_!(X_{\overline K},\mathbb Q_\ell)$, $X_K$ séparé de t.f. sur $K$.} \deleted{Pour} Soit $U$ un ouvert \deleted{affine} \inserted{dense} de $X_K$, $Y$ son complémentaire. Alors on a la suite exacte
\[
\cdots\longrightarrow H^i_!(\deleted{\unsure{U}}\unsure{Y_{\overline K}|U_{\overline K}},\deleted{\unsure{\mathbb Z_\ell|\mathbb Q_\ell}}\mathbb Q_\ell)\longrightarrow H^i_!(X_{\overline K},\mathbb Q_\ell)-H^i_!(U_{\overline K},\mathbb Q_\ell)-H^{i+1}_!(U_{\overline K},
\]
% The right edge of the sheet cuts the last displayed term after the comma.
qui montre, moyennant l'hyp. de récurrence, que
\scanpage{47}
\mspage{45}
% archive sheet number: 46
l'action de $I$ sur $H^i_!(X_{\overline K})$ est quasi unipotente ssi son action sur $H^i_!(U_{\overline K})$ l'est : caractère \uline{birationnel} de la question.

Si donc on peut trouver un \deleted{\unsure{modèle}} \inserted{``complété'' propre et lisse $\widehat U_K$} \deleted{\unsure{lisse}} de \deleted{\unsure{$X_K$}} $U_K$, alors on est ramené au cas d'un schéma \uline{lisse} et \uline{propre}. Si on admet la résolution des singularités \deleted{\unsure{par}} \inserted{(\unsure{par voie forte})}, sur la clôture \deleted{\unsure{alg.}} \inserted{parfaite} de $K$, pour les schémas de t.f. de dim $\leq n$, on aura gagné \inserted{(prenant $U$ affine)}, à condition que l'hypothèse de récurrence reste vérifiée après toute extension \inserted{radicale} finie \deleted{\unsure{ext}} $K'$ de $K$ (en prenant la clôture normale $V'$ de $V$ dans $K'$).

\inserted{\uline{Remarque} 9.3.} Donc en vertu de \deleted{\unsure{n°}} \unsure{7.1|9.1}, ça marche aussi pour $H^i(X_{\overline K},\unsure{\mathbb Q_\ell|Q_\ell})$, si $X$ lisse \uline{séparé}. Le cas non séparé s'y ramène, par suite spectrale de Leray d'un recouvrement fini de $X$ par des ouverts affines.

\deleted{9.3. Cas de $H^i_!(X_{\overline K},\mathbb Q_\ell)$ pour $X$ \inserted{séparé de t.f.} quelconque. \deleted{\unsure{Par ailleurs},} \deleted{On peut supposer $X$ réduit, et si} $U$ un ouvert dense, on \deleted{\unsure{peut}} \unsure{est} \unsure{comme} \deleted{\unsure{cela|alors}} \unsure{précédemment} \unsure{quitte pour|quitte à} remplacer $X$ par $U$. On peut trouver $U$ tel que \deleted{\unsure{changer $U$}} sur une \deleted{\unsure{une}} extension \inserted{finie} radicale $K'$ convenable de $K$, on ait $(\deleted{\unsure{X}}U_{K'})_{\mathrm{red}}$ \uline{lisse} sur $K'$. On est ramené au cas traité dans 9.2.}

\deleted{9.3. \uline{Cas de $H^i(X_{\overline K},\mathbb Q_\ell)$ pour $X$ de t.f. quelconque.}}
\scanpage{48}
\mspage{46}
% archive sheet number: 47
On trouve :

\uline{Proposition} 9.4. \inserted{a)} Pour \deleted{\unsure{établir}} la validité du th. de monodromie \inserted{pour les $H^i_!(X_{\overline K},\unsure{\mathbb Q_\ell|Q_\ell})$}, pour $V$ fixé, \deleted{et} des $X_K$ avec $\dim X_K\leq n$, et tout $i$, il suffit de savoir qu'il est vrai pour tout $X_{K'}$, où $X$ est lisse \inserted{de dim $\leq n$} sur \deleted{$\unsure{K'}$} \inserted{une} extension $K'$ finie radicale de $K$, considéré comme corps des fractions de la clôture normale $V'$ de $V$ dans $K'$. \deleted{\unsure{Et}} On peut \deleted{supposer que les $X_{K'}$} se borner à supposer de plus les $X_{K'}$ affines, \deleted{\unsure{et}} \inserted{b)} Dans le cas où \deleted{la résolution des singularités si} toute extension de t.f. de degré de transcendance $\leq n$ sur la clôture parfaite $K^{\unsure{p^{-\infty}|p^\infty}}$ de $K$ admet un modèle projectif et lisse, il suffit de supposer \inserted{les} $X_{K'}$ projectifs. c) Enfin, \deleted{si} les conditions \deleted{\unsure{précédentes}} \inserted{\deleted{\unsure{ci-dessus}}} de a) sont vérifiées, alors pour $\forall$ $X_K$ de t.f. et \uline{lisse} sur $K$, \inserted{de dim $\leq n$,} le th. de monodromie est vrai aussi pour les $H^i_{\deleted{!}}(X_{\overline K},\mathbb Q_\ell)$.

9.5. Nous allons maintenant examiner le th de monodromie pour $H^i(X_{\overline K},\mathbb Q_\ell)$, \deleted{\unsure{donc}} \inserted{pour} $X_K$ de t.f., en supposant que \deleted{les conditions de 9.4. a) sont satisfaites \unsure{pour tout $n$|et tout $n$}. On procède par récurrence sur $i$, le cas $i<0$ étant trivial,} c'est prouvé pour les \deleted{\unsure{les}} $H^{\unsure{j|i}}(X_{\overline K},\mathbb Q_\ell)$ pour $j<i$.
\scanpage{49}
\mspage{47}
% archive sheet number: 48
On peut supposer si on veut $\unsure{U|X}$ affine, donc quasi-projectif.

\deleted{\unsure{En}} On admet \deleted{que} la résolution des singularités sur la clôture parfaite \deleted{de $K$} de $K$ en dim $\leq n$. Alors, quitte à passer à une extension radicale finie de $K$, on pourra trouver \inserted{un} morphisme surjectif \deleted{et} propre \inserted{birationnel}
\[
X'\longrightarrow X,
\]
avec $X'$ lisse. \deleted{\unsure{et}} La suite spectrale de descente de Deligne nous donne
\[
H^{\unsure{n|i}}(X_{\overline K},\unsure{\mathbb Q_\ell|Q_\ell})\Longleftarrow E^{pq}_2=H^p(\alpha\longmapsto H^q((X'/X)^{\alpha+1}_{\overline K},\mathbb Q_\ell).
\]
\deleted{Dans} \deleted{Cette suite spectrale}
\deleted{Nous supposons que l'on connaît le th. de monodromie pour $H^j$ pour $j<i$, et sans restriction sur la dimension de $X$.} L'hyp. de récurrence nous assure que \deleted{$E^{pq}_2$ \unsure{satisf}} l'action de $I$ sur $E^{pq}_2$ est quasi-unipotente quand $q<i$. \deleted{Pour $q=i$, on doit prendre} Parmi les termes $E^{pq}_2$ avec $p+q=i$, le seul dont on ne sache encore \deleted{\unsure{qu'il est}} si l'action de $I$ dessus est quasi unipotente est donc $E^{0,i}_2\subset H^i(X'_{\overline K},\mathbb Q_\ell)$.

\deleted{Mais} On trouve donc

\uline{Prop.} 9.6. \deleted{\inserted{\unsure{Supposons rés. des sing. sur} $K^{\unsure{p^{-\infty}|p^\infty}}$ \unsure{des schémas de tte dimension|des schémas de t.f. dim $\leq n$}.} \deleted{Pour} Supposons le th. de monodromie pour les $H^i(X_{\overline K},\mathbb Q_\ell)$ soit vrai pour $i<n$, \inserted{pour $X_K$ de t.f.,} \deleted{il faut et suffit qu'il soit vrai} \deleted{Alors il est vrai pour $\forall$} \deleted{\unsure{pour $\forall$ $X_K$}} et pour $i=n$, quand $X_K$ est lisse \deleted{\unsure{séparé de t.f.}} \inserted{quasi-projectif}, et \unsure{ceci étant|reste-t-il} après extension \unsure{pure|purement} inséparable $K'$ de $K$. \deleted{\unsure{et pour lisse}} \deleted{Alors} Alors il est vrai pour $H^n(X_{\overline K},\mathbb Q_\ell)$ pour les $X$ propres qui, sur la clôture \deleted{\unsure{inséparable}, \unsure{de}} $K^{p^{-\infty}}$, admet une résolution des singularités.}
